\section{Validation and Box-Opening Strategy}
\label{sec:validation}

The signal region is blinded.
Control regions used for validation, none of which enters the fit:
% PROVENANCE
%   status      : ANALYSIS_DECISION
%   produced-by : utilities.py (Dst_veto_cut, lgb_comb, offline_cut);
%                 Recon_scripts/2_Reconstruction.py (wrong-charge sample)
%   commit      : dd520a5
%   sample      : MC16rd + Proc16/Prompt16, Run 1 + Run 2, e and mu
%   selection   : per row
%   corrections : n/a
\begin{table}[htbp]
  \centering
  \caption{Control regions and their intended tests.}
  \label{tab:control-regions}
  \small
  \begin{tabular}{p{0.34\textwidth} p{0.58\textwidth}}
    \hline
    Region & Purpose \\
    \hline
    $q^{2} < 3$ & $B \to D \ell \nu$ dominant; checks the core fit observables \\
    $D^{*}$ veto region, $0.135 < \Delta M < 0.145$ & $B \to D^{*} \ell \nu$ dominant; currently limited by the missing slow-$\pi^{0}$ correction \\
    $\Upsilon(4S)$ off-resonance & Continuum modelling \\
    BDT combinatorial-enhanced sideband & Generic-$B\bar{B}$ modelling \\
    Wrong-charge sample & Intended as an independent check of the tuned generic-$B\bar{B}$ weights; see the caveat in Sec.~\ref{sec:bbbar:validation} \\
    $D$-mass sidebands & Fake-$D$ modelling; also a fit channel (Sec.~\ref{sec:fit}) \\
    \hline
  \end{tabular}
\end{table}

\AnalysisTBD{prefit and postfit closure criteria, and the $e$ versus $\mu$
split-sample comparison}
\AnalysisTBD{box-opening plan and the criteria that must be satisfied before
unblinding; a staged (10\%) unblinding is under discussion}

\subsection{Known open issues}
\label{sec:validation:issues}
\AnalysisTBD{$E_{\mathrm{cms}}$ and experiment-level data/MC discrepancies in
several Run~1 and Run~2 experiments, hypothesised to originate from
background-overlay run mismatches; and the fixed, non-run-dependent
$B\bar{B}$ cross section in MC16rd despite the $E_{\mathrm{cms}}$ drift.
Sources to be recorded}
